\section*{अध्याय \ref{ch:b1:gas-exchange} --- श्वसनी गैस विनिमय}

\begin{solution}{exo:b1:gas-exchange:1}
$\dot V = K(A/x)(P_1 - P_2)$: बड़ा क्षेत्रफल, पतली रोक,
ऊँचा बाहरी \omterm{def:b1:gas-exchange:partial}{आंशिक दाब} (संवातन) और नीचा भीतरी दाब
(छिड़काव)।
\end{solution}

\begin{solution}{exo:b1:gas-exchange:2}
प्रति लीटर तीस गुना कम ऑक्सीजन (\qty{210}{mL/L} के सामने \qty{7}{});
आठ सौ गुना सघन और पचास गुना अधिक श्यान, इसलिए पंप करना महँगा;
और ऑक्सीजन उसमें दस हज़ार गुना धीमे विसरित होती है, इसलिए
माध्यम को पृष्ठ के ठीक आरपार चलाना ही पड़ता है।
\end{solution}

\begin{solution}{exo:b1:gas-exchange:3}
सहधारा: जल और रक्त दोनों \qty{12}{kPa} पर निकलते हैं; निष्कर्षण
$(21 - 12)/21 = \qty{43}{\%}$। \omterm{def:b1:gas-exchange:gill}{प्रतिधारा}: जल \qty{4}{kPa} पर निकलता है,
रक्त \qty{18}{kPa} पर; निष्कर्षण $(21 - 4)/21 =
\qty{81}{\%}$।
\end{solution}

\begin{solution}{exo:b1:gas-exchange:4}
घुली हुई (\qty{7}{\%}), हीमोग्लोबिन के ऐमीनो समूहों से बँधी
(\qty{23}{\%}), और प्लाज़्मा में बाइकार्बोनेट के रूप में (\qty{70}{\%})।
\end{solution}

\begin{solution}{exo:b1:gas-exchange:5}
मिनट-संवातन \qty{6.75}{L/min}; कूपिकीय $15\times 0.3 = \qty{4.5}{L/min}$;
ग्रहण $4.5\times 0.07 = \qty{0.32}{L/min}$। उथली: मिनट-संवातन तब भी
\qty{6.75}{L/min}, पर कूपिकीय $30\times 0.075 =
\qty{2.25}{L/min}$ और ग्रहण \qty{0.16}{L/min} --- यानी उसी
श्वास-श्रम के लिए आधा; मृत स्थान की कीमत हर साँस पर चुकानी पड़ती है।
\end{solution}

\begin{solution}{exo:b1:gas-exchange:6}
प्रति घंटा $2/(6\times 0.75) = \qty{0.44}{L}$, यानी \qty{20}{mL} के
अपने आयतन का बाईस गुना।
\end{solution}

\begin{solution}{exo:b1:gas-exchange:7}
$A/x$ $2\times 4 = 8$ से गिरता है: \qty{31}{mL/min}। \qty{250}{} बहाल
करने के लिए अंतर \qty{64}{kPa} होना चाहिए --- जो असंभव है, क्योंकि
कूपिकीय दाब हवा में \qty{13}{kPa} से ऊपर नहीं जा सकता और शिरा-रक्त
शून्य से नीचे नहीं गिर सकता; रोगी उसे केवल ऑक्सीजन साँस लेकर चढ़ा सकता
है, और तब भी आठ गुना नहीं।
\end{solution}

\begin{solution}{exo:b1:gas-exchange:8}
\omterm{prop:b1:gas-exchange:transport}{कार्बोनिक ऐनहाइड्रेज़} \omterm{def:b1:cell-unit-of-life:cell}{कोशिका} के भीतर है, इसलिए बाइकार्बोनेट वहीं बनता है
और जमा होता है; वह अपनी प्रवणता के नीचे किसी ऐसे विनिमायक से होकर बाहर
विसरित होता है जो हर बाइकार्बोनेट के बदले एक क्लोराइड भीतर लेता है, और
आवेश संतुलित रखता है। उस विनिमय के बिना बाइकार्बोनेट अपने प्रोटॉन के
साथी के साथ भीतर ही रह जाता, \omterm{def:b1:cell-unit-of-life:cell}{कोशिका} की \omterm{def:b1:membranes-transport:osmosis}{परासरणिता} चढ़ा देता, और जल
भीतर घुसता: \omterm{def:b1:cell-unit-of-life:cell}{कोशिका} \omterm{def:b1:body-plans-tissues:tissue}{ऊतकों} में फूल जाती (वह वैसे भी ज़रा फूलती है)।
\end{solution}

\begin{solution}{exo:b1:gas-exchange:9}
साँस लेने की तलब चढ़ती $\mathrm{CO_2}$ से आती है। अतिसंवातन
$\mathrm{CO_2}$ को सामान्य से बहुत नीचे ले जाता है और ऑक्सीजन ज़्यादा
नहीं जोड़ता (हीमोग्लोबिन पहले ही संतृप्त था); डुबकी के दौरान ऑक्सीजन
अचेतना के स्तर तक गिर जाती है इससे पहले कि $\mathrm{CO_2}$ उस स्तर तक
चढ़ आए जो साँस लेने पर मजबूर करता है। बिना अतिसंवातन के
$\mathrm{CO_2}$ उस स्तर तक तब पहुँचती है जब ऑक्सीजन अब भी भरपूर है, और
गोताखोर पानी के ऊपर आ जाता है।
\end{solution}

\begin{solution}{exo:b1:gas-exchange:10}
$A = \pi(10^{-5})^2 = \qty{3.1e-10}{m^2}$; $\Delta c = JL/DA =
10^{-9}\times 10^{-3}/(2\times 10^{-5}\times 3.1\times 10^{-10}) =
\qty{160}{mol/m^3}$ --- यानी उपलब्ध \qty{8.7}{} से कहीं अधिक, इसलिए
इतनी सँकरी नली \qty{1}{mm} पर पेशी को विसरण से नहीं पाल सकती;
कोई चौड़ी नली (\qty{100}{\micro m}: $\Delta c = \qty{6.4}{mol/m^3}$, बस
किसी तरह संभव) या सक्रिय संवातन चाहिए। \qty{10}{mm} पर
यह माँग फिर से दस गुनी है। चूँकि नली की लंबाई शरीर के आकार के साथ बढ़ती
है और माँग आयतन के साथ, श्वासनलिकाओं से होकर विसरण कीटों के आकार पर कुछ
ही सेंटीमीटर की सीमा बाँध देता है, और उससे अधिक केवल पंपन से।
\end{solution}

\begin{solution}{exo:b1:gas-exchange:11}
मछली: जल एक दिशा में, रक्त उसके विपरीत (\omterm{def:b1:gas-exchange:gill}{प्रतिधारा}),
निष्कर्षण \qty{80}{\%}। पक्षी: हवा एक दिशा में परानलिकाओं से होकर,
रक्त समकोण पर (अनुप्रस्थ धारा), निष्कर्षण
मध्यवर्ती पर स्तनधारियों से ऊँचा, और विनिमय-पृष्ठ को
बासी हवा कभी नहीं मिलती। कोई ज्वारीय \omterm{def:b1:gas-exchange:lung}{फेफड़ा} ताज़ी हवा को अवशिष्ट गैस से
मिला देता है, इसलिए विनिमय-पृष्ठ बहुत हुआ तो कूपिकीय मिश्रण देखता है, और
रक्त उसी से साम्य में आ सकता है: पृष्ठ चाहे जो हो, धमनीय रक्त कूपिकीय
दाब से आगे नहीं जा सकता, और जब तक कोई अवशिष्ट आयतन बचा है तब तक कूपिकीय
दाब साँस में ली गई हवा के दाब के पास नहीं आ सकता।
\end{solution}

\begin{solution}{exo:b1:gas-exchange:12}
$\mathrm{CO_2}$ बेहतर संकेत है क्योंकि उसका धमनीय स्तर संवातन पर सीधे
और तीखे ढंग से जवाब देता है (संवातन आधा कीजिए और वह दुगना हो जाता
है), वह केंद्रीय ग्राहियों द्वारा pH के रास्ते परिशुद्धि से नापा जाता है,
और ऑक्सीजन, जो हीमोग्लोबिन के वक्र की चपटी चोटी पर बैठी है, संवातन के
सामान्य परास भर कम ही बदलती है: कोई $\mathrm{CO_2}$ तापस्थापी ऑक्सीजन को
उप-उत्पाद के रूप में ठीक रखता है। वह तब विफल होता है जब ऑक्सीजन
$\mathrm{CO_2}$ से स्वतंत्र होकर गिरे --- ऊँचाई पर,
जहाँ साँस में ली ऑक्सीजन नीची है पर $\mathrm{CO_2}$ का उत्पादन नहीं,
इसलिए $\mathrm{CO_2}$ का संकेत कहता है ``कम साँस लो'' जबकि ऑक्सीजन कहती है
``अधिक साँस लो''; और केवल ग्रीवा-पिंड, जब ऑक्सीजन बहुत नीची हो, उस पर
हावी होते हैं, और उस समझौते (क्षारमयता के साथ अतिसंवातन) को बैठने में
वृक्कों के समायोजन के दिन लगते हैं।
\end{solution}

\begin{solution}{pb:b1:gas-exchange:1}
\textbf{1.} हवा $210/22.4 = \qty{9.4}{mmol/L}$; जल $7/22.4 =
\qty{0.31}{mmol/L}$।
\textbf{2.} 30.
\textbf{3.} हवा: $1.2/9.4 = \qty{0.13}{g}$; जल: $1000/0.31 =
\qty{3200}{g}$ --- यानी \num{25000} गुना अधिक द्रव्यमान।
\textbf{4.} जल की $210/7 = \qty{30}{L}$।
\textbf{5.} जल इतना भारी और ऑक्सीजन में इतना गरीब है कि उसे भीतर-बाहर
ठेला नहीं जा सकता; उसे एक ही बार गुज़ारना पड़ता है, जो \omterm{def:b1:gas-exchange:gill}{प्रतिधारा} को भी
संभव बनाता है; और हवा इतनी सस्ती है कि दो बार चलाई जा सके।
\textbf{6.} मिनट-संवातन $12\times 0.5 = \qty{6}{L/min}$; कूपिकीय
$12\times 0.35 = \qty{4.2}{L/min}$।
\textbf{7.} हर साँस का, और मिनट-संवातन का, $150/500 = \qty{30}{\%}$।
\textbf{8.} पहुँचाई $4.2\times 0.21 = \qty{0.88}{L/min}$; ग्रहण
$4.2\times 0.07 = \qty{0.29}{L/min}$, जो \qty{250}{mL/min} के पास है।
\textbf{9.} $0.25/(6\times 0.21) = \qty{20}{\%}$ (\omterm{def:b1:gas-exchange:lung}{कूपिकाओं} तक पहुँचने
वाली का \qty{28}{\%})।
\textbf{10.} $0.14\times(101 - 6.3) = \qty{13.3}{kPa}$।
\textbf{11.} $250/5 = \qty{50}{mL/L}$; शिरा-रक्त \qty{150}{mL/L},
यानी \qty{75}{\%} संतृप्त।
\textbf{12.} $3000/20 = \qty{150}{mL/L}$; शिरा-रक्त \qty{50}{mL/L},
यानी \qty{25}{\%} संतृप्त।
\textbf{13.} $50/60 = \qty{0.83}{mL/min}$।
\textbf{14.} $0.83/(7\times 0.8) = \qty{0.149}{L/min}$, यानी प्रति घंटा
\qty{8.9}{L}।
\textbf{15.} $0.83/(7\times 0.43) = \qty{0.28}{L/min}$: \omterm{def:b1:gas-exchange:gill}{प्रतिधारा}
पंप किए जाने वाले जल को, और श्रम को, आधा कर देती है।
\textbf{16.} प्रति घंटा \qty{8.9}{kg} जल; और मनुष्य प्रति घंटा $6\times
60\times 1.2 = \qty{430}{g}$ हवा चलाता है --- यानी द्रव्यमान का बीसवाँ
भाग, और तीन सौ गुनी ऑक्सीजन के लिए (प्रति घंटा \qty{50}{mL} के सामने
$15\,000$)।
\textbf{17.} $100\times(0.95 - 0.15) = \qty{80}{mL/L}$; हृद-निर्गम
$0.83/0.080 = \qty{10.4}{mL/min}$।
\textbf{18.} जल/रक्त $149/10.4 = 14$; हवा/रक्त $4.2/5 = 0.84$।
मछली को प्रति आयतन रक्त जल के चौदह आयतन गुज़ारने ही पड़ते हैं
क्योंकि हर एक इतनी कम ऑक्सीजन रखता है; और मनुष्य को लगभग एक।
\textbf{19.} ग्रहण $\times 1.5$, और जल $5.5/7 = 0.79$ गुना ही रखता है:
इसलिए संवातन $\times 1.9$; और साँस लेने की लागत, जो मोटे तौर पर
बहाव के अनुपात में है (या अधिक, क्योंकि वह रैखिक से तेज़ चढ़ती है),
किसी बड़े ग्रहण के \qty{10}{\%} से कम से कम \qty{13}{\%} पर चली जाती है
--- यानी मछली गरम पानी में साँस लेने के लिए अधिक मेहनत करती है।
\textbf{20.} जल के प्रवेश वाले सिरे पर: $21 - 18 = \qty{3}{kPa}$; जल के
निकास वाले सिरे पर: $4 - 3 = \qty{1}{kPa}$: दोनों सिरों पर धनात्मक।
\textbf{21.} प्रवेश वाले सिरे पर $21 - 3 = \qty{18}{kPa}$; निकास वाले
सिरे पर $12 - 12 =
0$: दोनों तरल साम्य में आ चुके हैं और कोई प्रवणता
नहीं बची, यद्यपि जल तब भी अपनी ऑक्सीजन का \qty{57}{\%} रखे है।
\textbf{22.} सहधारा: रक्त किसी एक ही, अच्छी तरह मिली कूपिकीय गैस से
साम्य में आ जाता है और उसी के दाब पर निकलता है; वह
\qty{13.3}{kPa} से आगे नहीं जा सकता क्योंकि आगे मिलने को कोई ताज़ी गैस
नहीं है, जैसी \omterm{def:b1:gas-exchange:gill}{प्रतिधारा} में होती है।
\textbf{23.} ट्राउट $2000/(2\times 10^{-4}) = \num{1e7}$ (प्रति cm cm$^2$);
मनुष्य $10^6/(5\times 10^{-5}) = \num{2e10}$: अनुपात \num{2000}। और ग्रहण
$250/0.83 = 300$। ट्राउट $A/x$ की प्रति इकाई सात गुना अधिक ऑक्सीजन लेती
है: इसलिए पटलिका के आरपार उसका माध्य दाब-अंतर बड़ा होना ही चाहिए
--- जो \omterm{def:b1:gas-exchange:gill}{प्रतिधारा} देती है, और समूचे पृष्ठ भर एक प्रवणता बनाए रखती है
जबकि ज्वारीय फेफड़े की प्रवणता साम्य के अंत के पास
छोटी रह जाती है।
\textbf{24.} $\mathrm{CO_2}$ जल में ऑक्सीजन से पच्चीस गुना अधिक
विलेय है: जल उसे उतनी ही तेज़ी से बहा ले जाता है जितनी तेज़ी से वह
बनती है, और किसी मछली के रक्त की $\mathrm{CO_2}$ बहुत नीची रहती है।
\textbf{25.} \omterm{def:b1:gas-exchange:gill}{क्लोम} साँस में लिए माध्यम की ऑक्सीजन का लगभग \qty{80}{\%},
\omterm{def:b1:gas-exchange:lung}{फेफड़ा} लगभग \qty{20}{\%}; और यह अंतर \omterm{def:b1:gas-exchange:gill}{क्लोम} के एकतरफ़ा बहाव तथा जल और
रक्त की उसकी \omterm{def:b1:gas-exchange:gill}{प्रतिधारा} व्यवस्था से पड़ता है।
\end{solution}
